\documentclass[11pt]{amsart}

\usepackage[T1]{fontenc}
\usepackage{lmodern}
\usepackage{microtype}
\usepackage{mathtools,amssymb,amsthm}
\usepackage[hidelinks]{hyperref}

\hypersetup{
  pdftitle={Alpha-Valuations of C4-Bouquets},
  pdfsubject={Graph labeling},
  pdfkeywords={graceful labeling, alpha-valuation, cycle amalgamation, bouquet}
}

\newtheorem{theorem}{Theorem}
\newtheorem{corollary}[theorem]{Corollary}
\DeclareMathOperator{\Amal}{Amal}

\title[$\alpha$-valuations of $C_4$-bouquets]
  {$\alpha$-Valuations of $C_4$-Bouquets}
\date{}
\subjclass[2020]{05C78}
\keywords{graceful labeling, $\alpha$-valuation, cycle amalgamation, bouquet}

\begin{document}

\begin{abstract}
For $n\ge1$, let $B_n$ be the bouquet obtained by identifying one vertex
in each of $n$ copies of $C_4$. We prove that $B_n$ admits an
$\alpha$-valuation if and only if $n\le4$. The existence half is known:
for $2\le n\le4$ it is the case $m=4$ of Theorem~3.8 of
Figueroa-Centeno, Ichishima, and Muntaner-Batle, and $B_1=C_4$ is an
$\alpha$-graph by Rosa's classification of cycles; what is new here is
that no $n\ge5$ admits one. In particular, $B_5$, with $16$ vertices and
$20$ edges, is a bipartite counterexample to the sufficiency direction of
Conjecture~3.10 of the same paper. Thus the conjecture fails even for
$m=4$.
\end{abstract}

\maketitle

\section{The result}

A graceful labeling of a graph with $q$ edges is an injection
$f\colon V(G)\to\{0,1,\dots,q\}$ whose induced edge differences are
$1,2,\dots,q$. It is an $\alpha$-valuation if there is an integer $k$
such that every edge has one endpoint labeled at most $k$ and the other
labeled greater than $k$.

Figueroa-Centeno, Ichishima, and Muntaner-Batle conjectured that the
one-point union $\Amal(C_m,v,n)$ admits an $\alpha$-valuation if and only
if $mn\equiv0\pmod4$ \cite[Conjecture~3.10]{FCIM}. The twenty-eighth
edition of Gallian's survey records the same conjecture
\cite[p.~53]{Gallian}. Because every graph admitting an
$\alpha$-valuation is bipartite, odd $m$ with $4\mid n$ are immediate
exceptions to the statement as printed. The nontrivial question is
whether the divisibility condition is sufficient for even $m$. The
theorem below answers this negatively already for $m=4$ and determines
all $n$ in that case.

The positive part of that determination is due to the authors of the
conjecture: if $m\equiv0\pmod4$ and $n\in\{2,3,4\}$, then
$\Amal(C_m,v,n)$ admits an $\alpha$-valuation
\cite[Theorem~3.8]{FCIM}. (Theorem~3.8 is stated there for
$n=1,2,3,4$, although $\Amal(G,v,n)$ is defined in \cite{FCIM} only for
$n\ge2$; the remaining case is $\Amal(C_4,v,1)=C_4$, an $\alpha$-graph
because $4\equiv0\pmod4$ \cite{Rosa}.) What is new below is the
nonexistence for every $n\ge5$.

For $n\ge1$, write
\[
  B_n=\Amal(C_4,v,n),
\]
and denote its common vertex by $h$. Label the remaining vertices so
that its $i$th square is
\[
  h x_i y_i z_i h,\qquad 1\le i\le n.
\]

\begin{theorem}\label{thm:main}
The graph $B_n$ admits an $\alpha$-valuation if and only if $n\le4$.
\end{theorem}

\begin{proof}
For $1\le n\le4$ an $\alpha$-valuation exists by
\cite[Theorem~3.8]{FCIM} (case $m=4$, $n\ge2$) and by \cite{Rosa} (case
$n=1$). The following table records explicit witnesses, so that the
argument here is self-contained and the positive direction can be checked
by hand. The triples are $(f(x_i),f(y_i),f(z_i))$, in increasing order of
$i$, and $k$ is the characteristic.
\[
\begin{array}{c|c|c|l}
 n & f(h) & k & (f(x_i),f(y_i),f(z_i))_{i=1}^{n} \\
\hline
 1 & 0 & 2 & (4,2,3) \\
 2 & 0 & 2 & (3,2,4),\ (6,1,8) \\
 3 & 1 & 3 & (6,0,12),\ (4,3,5),\ (9,2,11) \\
 4 & 2 & 4 & (7,1,9),\ (15,0,16),\ (5,4,6),\ (12,3,14)
\end{array}
\]
In each row the labels are injective, the two partite sets are separated
at $k$, and the edge differences are $1,2,\dots,4n$.

The remainder of the proof is the other direction. Suppose that $B_n$ has
an $\alpha$-valuation $f$. Its bipartition is
\[
  A=\{h,y_1,\dots,y_n\},
  \qquad
  B=\{x_i,z_i:1\le i\le n\}.
\]
Since $B_n$ is connected, the lower and upper label classes are its two
partite sets. Replacing $f$ by $4n-f$ if necessary, assume that $A$ is
the lower class. Set
\[
  S=f(A),\qquad T=f(B),\qquad a=f(h),\qquad U=S\setminus\{a\}.
\]
Thus every element of $S$ is smaller than every element of $T$.
For $t\in T$, let $g(t)$ be the label of the vertex $y_i$ in the same
square as the vertex of $B$ labeled $t$. Then
\[
  g^{-1}(f(y_i))=\{f(x_i),f(z_i)\},
\]
so every fiber of $g\colon T\to U$ has size two.

The edge differences decompose as
\begin{equation}\label{eq:partition}
  \Delta_h=\{t-a:t\in T\},
  \qquad
  \Delta_y=\{t-g(t):t\in T\},
  \qquad
  \Delta_h\mathbin{\dot\cup}\Delta_y=\{1,\dots,4n\}.
\end{equation}
The difference $4n$ forces the labels $0$ and $4n$ to occur. Hence
$0\in S$ and $4n\in T$.

We first show that
\begin{equation}\label{eq:a}
  a\le2.
\end{equation}
This is immediate if $a=0$. If $a>0$, then
$\max\Delta_h=4n-a$, so every element of
\[
  I=\{4n-a+1,\dots,4n\}
\]
lies in $\Delta_y$. For each $d\in I$, let $t_d$ be the unique $t\in T$ with $d=t-g(t)$; it
is unique because the differences in \eqref{eq:partition} are pairwise
distinct. We call $t_d$ the label \emph{serving} $d$. Since
$t_d=d+g(t_d)$, we have $t_d\in I$. The labels
$t_d$ are distinct, and $I$ has $a$ elements; therefore
$\{t_d:d\in I\}=I$. Consequently,
\[
  \sum_{d\in I}g(t_d)
  =\sum_{d\in I}t_d-\sum_{d\in I}d=0.
\]
Thus every $g(t_d)$ equals $0$, so $0\in U$. Since
$|g^{-1}(0)|=2$, we obtain $a=|I|\le2$.

Let
\[
  M=\max S,
  \qquad
  \tau=\min T.
\]
We next show that
\begin{equation}\label{eq:span}
  M-a\le2.
\end{equation}
Since $\min\Delta_h=\tau-a$, every integer
$1,\dots,\tau-a-1$ lies in $\Delta_y$. If
$d=t_d-g(t_d)$ is one of these differences, then
\[
  \tau\le t_d=d+g(t_d)\le\tau-a-1+M.
\]
The serving labels $t_d$ are distinct, while this interval contains
$M-a$ integers. Hence
\[
  \tau-a-1\le M-a
\]
(if $\tau-a-1\le0$ there are no such differences and the inequality is
trivial, since $a\in S$ gives $M\ge a$). Because $\tau>M$, it follows
that $\tau=M+1$.

Put $r=M-a$. If $r=0$, \eqref{eq:span} is immediate. If $r>0$, the
small differences are $1,\dots,r$, and their serving labels exhaust
$M+1,\dots,M+r$. Therefore
\[
  \sum_{d=1}^{r}g(t_d)
  =\sum_{j=1}^{r}(M+j)-\sum_{d=1}^{r}d
  =rM.
\]
Since $r>0$, we have $M\in U$. Every summand is at most $M$, so all
of them equal $M$. Hence $r\le|g^{-1}(M)|=2$, proving
\eqref{eq:span}.

By \eqref{eq:a} and \eqref{eq:span}, $M\le4$. The set $S$ consists of
$n+1$ distinct labels in $\{0,1,\dots,M\}$, so $n\le M\le4$.
\end{proof}

Both inequalities are attained by the $n=4$ row of the table, so the
proof is self-certifying at its own boundary: there $a=2$, $M=4$,
$\tau=M+1=5$, and the two fibers named in the proof are
$g^{-1}(0)=\{15,16\}$ and $g^{-1}(4)=\{5,6\}$. Thus \eqref{eq:a} and
\eqref{eq:span} are simultaneously saturated by an exhibited labeling,
and the chain $n\le M=a+(M-a)\le2+2$ cannot be improved at either step.

\begin{corollary}\label{cor:counterexample}
For every $n\ge5$, the bipartite graph $B_n$ has $4n$ edges and admits no
$\alpha$-valuation. Hence $B_5$ refutes the conjectured sufficiency
condition within the even-cycle subfamily of Conjecture~3.10 in
\cite{FCIM}.
\end{corollary}

\begin{proof}
For $m=4$, the conjectured divisibility condition holds for every $n$,
whereas Theorem~\ref{thm:main} excludes every $n\ge5$.
\end{proof}

\subsection*{Status}

What is refuted is Conjecture~3.10 of \cite{FCIM} as stated, and what is
in addition settled is its $m=4$ column, completely. The existence half
of Theorem~\ref{thm:main} is not ours: it is \cite[Theorem~3.8]{FCIM} for
$n=2,3,4$ and \cite{Rosa} for $n=1$, and the table above only exhibits
witnesses for it. Nothing else is claimed. The nonexistence argument uses
$m=4$ essentially --- each vertex of $B$ is adjacent only to $h$ and to
the $y$ of its own square, which fails for every longer cycle --- so the
columns $m\equiv0\pmod4$ with $m\ge8$, and $m\equiv2\pmod4$ with $n$ even
and $n\ge6$, are untouched here. In the former columns the cases $n\le4$
are already settled positively by \cite[Theorem~3.8]{FCIM}, so what
remains open there is the range $n\ge5$. No minimality claim is
made: $B_5$ is exhibited as a bipartite counterexample, not asserted to be
the smallest counterexample of any kind.

The obstruction is specific to the $\alpha$-property and not to
gracefulness: $\Amal(C_4,v,n)$ is graceful for every $n$, by a result of
Shee recorded in \cite[p.~19]{Gallian}. So the conjecture of Koh, Rogers,
Lee, and Toh that $\Amal(C_m,v,n)$ is graceful if and only if
$mn\equiv0$ or $3\pmod4$, from which Conjecture~3.10 is descended, is not
affected by anything above.

\begingroup
\small
\raggedright
\begin{thebibliography}{9}

\bibitem{FCIM}
R.~M. Figueroa-Centeno, R.~Ichishima, and F.~A. Muntaner-Batle,
\emph{Labeling the vertex amalgamation of graphs},
Discuss. Math. Graph Theory \textbf{23} (2003), 129--139.
\href{https://doi.org/10.7151/dmgt.1190}{doi:10.7151/dmgt.1190}.

\bibitem{Gallian}
J.~A. Gallian,
\emph{A dynamic survey of graph labeling},
Electron. J. Combin. (2025), Dynamic Survey \#DS6,
twenty-eighth edition, October 30, 2025.
\href{https://doi.org/10.37236/27}{doi:10.37236/27}.

\bibitem{Rosa}
A.~Rosa,
\emph{On certain valuations of the vertices of a graph},
Theory of Graphs (Internat.\ Sympos., Rome, 1966),
Gordon and Breach, New York; Dunod, Paris, 1967, pp.~349--355.

\end{thebibliography}
\endgroup

\end{document}